\clearpage
\phantomsection
\addcontentsline{toc}{chapter}{Simulado e Revisão Integrada do Módulo 4}
\chapter*{Simulado e Revisão Integrada: Módulo 4}
\markboth{SIMULADO DO MÓDULO 4}{SIMULADO DO MÓDULO 4}

\noindent\tagCampo\ \tagPratica
\vspace{6pt}

\begin{boxConceito}[Instruções do Simulado Integrado (Capítulos 10, 11 e 12)]
Este simulado integra os conteúdos das Semanas 11 a 13 (Indicadores Ambientais e PEIR no Semiárido Baiano, Qualidade do Solo e Diagnóstico de Erosão Linear/USLE), compondo a base técnica direta do seu Trabalho Prático Final.
\begin{itemize}[leftmargin=1.5em, itemsep=2pt]
  \item \textbf{Tempo sugerido:} $90$ minutos sem consulta ao gabarito.
  \item \textbf{Estrutura:} $8$ questões objetivas autorais no estilo ENADE e concursos ambientais (valor: $0{,}75$ ponto cada $= 6{,}0$ pontos) e $2$ questões dissertativas aplicadas (valor: $2{,}0$ pontos cada $= 4{,}0$ pontos).
  \item Ao concluir, confira o \textbf{Gabarito Comentado}, o \textbf{Espelho Dissertativo} e preencha a \textbf{Tabela de Diagnóstico de Revisão} na página seguinte.
\end{itemize}
\end{boxConceito}

\section*{Parte I: Questões Objetivas (Estilo ENADE e Concursos Ambientais)}

\begin{enumerate}[leftmargin=1.8em, itemsep=8pt]
  \item \textbf{[Cap. 10, Critérios de Indicadores]} Na estruturação de um plano de monitoramento ambiental conforme OECD (2003), Fidalgo (2003) e IBGE (2018), um técnico propôs substituir um protocolo laboratorial de altíssimo custo e baixa periodicidade por medições semestrais padronizadas de cobertura vegetal e densidade de focos erosivos por quilômetro de estrada. Essa substituição fortalece principalmente qual critério de seleção de indicadores?
  \begin{alternativas}
    \item Subjetividade interpretativa.
    \item Mensurabilidade e viabilidade operacional de acompanhamento temporal.
    \item Esteticismo paisagístico.
    \item Independência frente às metas de gestão.
    \item Exclusividade para áreas urbanas impermeabilizadas.
  \end{alternativas}

  \item \textbf{[Cap. 10, Modelo PEIR e Campo no Semiárido]} No diagnóstico integrado de um complexo eólico na região de Canudos (BA), associaram-se corretamente os quatro compartimentos do modelo PEIR. Assinale a alternativa em que há correspondência técnica exata:
  \begin{alternativas}
    \item Pressão: colisão de aves; Estado: PRAD executado; Impacto: velocidade dos ventos alísios; Resposta: supressão vegetal.
    \item Pressão: terraplenagem de platôs e abertura de acessos na Caatinga; Estado: aumento da fração de solo exposto e da densidade do solo; Impacto: erosão linear em taludes e fragmentação de rotas da avifauna; Resposta: drenagem pluvial, bioengenharia e zonas de exclusão para aves ameaçadas.
    \item Pressão: zonas de exclusão para avifauna; Estado: ravinas profundas; Impacto: abertura de acessos; Resposta: solo exposto.
    \item Pressão: chuvas convectivas naturais; Estado: PRAD; Impacto: densidade viária; Resposta: assoreamento de barreiros.
    \item Pressão: geração de energia elétrica; Estado: conflitos fundiários; Impacto: cobertura vegetal nativa; Resposta: terraplenagem de novos platôs.
  \end{alternativas}

  \item \textbf{[Cap. 10, Eixo Energético do Submédio São Francisco]} Em relação ao eixo territorial que compreende Sobradinho, Casa Nova, Sento Sé, Juazeiro, Rodelas e Paulo Afonso, é correto afirmar que:
  \begin{alternativas}
    \item A expansão eólica nessa faixa ocorreu de forma desvinculada da rede de transmissão construída historicamente a partir das usinas hidrelétricas do São Francisco.
    \item Os impactos hidrelétricos restringem-se ao local exato das barragens, sem afetar o pulso natural de cheias, a agricultura de vazante ou os usos múltiplos da água a jusante.
    \item A convivência entre matriz hidrelétrica, parques eólicos, perímetros irrigados e comunidades locais demanda indicadores ambientais integrados que articulem gestão hídrica, conservação da Caatinga e ordenamento territorial.
    \item O bioma predominante nesse eixo é a Floresta Ombrófila Densa, o que impede a ocorrência de déficits hídricos sazonais.
    \item A regularização de vazões elimina todos os conflitos entre geração de energia e agricultura irrigada durante estiagens prolongadas.
  \end{alternativas}

  \item \textbf{[Cap. 11, Física e Química do Solo]} Uma amostra indeformada coletada em anel de $100\text{ cm}^3$ apresentou massa seca de $159\text{ g}$ (adote $D_p = 2{,}65\text{ g}\cdot\text{cm}^{-3}$), e a análise química indicou teor de matéria orgânica $MO = 17{,}24\text{ g}\cdot\text{kg}^{-1}$. Os valores de porosidade total ($P_t$) e de carbono orgânico total ($CO$, pelo fator de van Bemmelen $1{,}724$) são, respectivamente:
  \begin{alternativas}
    \item $30{,}0\%$ e $10{,}0\text{ g}\cdot\text{kg}^{-1}$.
    \item $40{,}0\%$ e $10{,}0\text{ g}\cdot\text{kg}^{-1}$.
    \item $40{,}0\%$ e $29{,}7\text{ g}\cdot\text{kg}^{-1}$.
    \item $60{,}0\%$ e $10{,}0\text{ g}\cdot\text{kg}^{-1}$.
    \item $60{,}0\%$ e $29{,}7\text{ g}\cdot\text{kg}^{-1}$.
  \end{alternativas}

  \item \textbf{[Cap. 11, Normalização e IQS]} Em uma área sob recuperação ambiental, mediram-se quatro indicadores de qualidade do solo comparados a uma referência local:
  \begin{itemize}[leftmargin=1.4em, itemsep=1.5pt]
    \item $MO$ (maior melhor): $24\text{ g}\cdot\text{kg}^{-1}$ (referência: $30\text{ g}\cdot\text{kg}^{-1}$);
    \item $TIB$ (maior melhor): $30\text{ mm}\cdot\text{h}^{-1}$ (referência: $60\text{ mm}\cdot\text{h}^{-1}$);
    \item $D_s$ (menor melhor): $1{,}50\text{ g}\cdot\text{cm}^{-3}$ (referência: $1{,}20\text{ g}\cdot\text{cm}^{-3}$);
    \item $DMP$ de agregados (maior melhor): $1{,}50\text{ mm}$ (referência: $3{,}00\text{ mm}$).
  \end{itemize}
  O Índice de Qualidade do Solo simples ($\text{IQS}_{\text{simples}}$) dessa área é igual a:
  \begin{alternativas}
    \item $0{,}500$.
    \item $0{,}600$.
    \item $0{,}650$.
    \item $0{,}725$.
    \item $0{,}800$.
  \end{alternativas}

  \item \textbf{[Cap. 12, Processos Erosivos Lineares]} Durante vistoria em uma encosta degradada, o geógrafo registrou: (i) crosta superficial fina e impermeabilizada nas áreas de solo nu; (ii) pequenos canais rasos de $0{,}25\text{ m}$ de profundidade ao longo de trilhas; e (iii) no terço inferior, uma grande fenda de $3{,}8\text{ m}$ de profundidade com surgência contínua de água subterrânea e blocos de solo desabando nas paredes. Os três fenômenos observados correspondem, respectivamente, a:
  \begin{alternativas}
    \item \textit{Piping}, voçoroca e erosão laminar.
    \item Selamento superficial, sulcos erosivos e voçoroca ativa.
    \item Assoreamento fluvial, ravina e salpicamento.
    \item Compactação profunda, voçoroca e sulco efêmero.
    \item Selamento superficial, voçoroca e ravina estável.
  \end{alternativas}

  \item \textbf{[Cap. 12, Equação Universal de Perda de Solo]} Uma encosta cultivada possui $R = 5.500\text{ MJ}\cdot\text{mm}\cdot\text{ha}^{-1}\cdot\text{h}^{-1}\cdot\text{ano}^{-1}$, $K = 0{,}02\text{ t}\cdot\text{ha}\cdot\text{h}\cdot\text{ha}^{-1}\cdot\text{MJ}^{-1}\cdot\text{mm}^{-1}$ e $LS = 2{,}0$. Atualmente, sob manejo convencional morro abaixo ($C = 0{,}30$; $P = 1{,}00$), a área sofre forte degradação. Caso seja implantado sistema conservacionista com cobertura permanente e plantio em contorno ($C = 0{,}05$; $P = 0{,}50$), a perda de solo anual passará de:
  \begin{alternativas}
    \item $66{,}0$ para $5{,}5\text{ t}\cdot\text{ha}^{-1}\cdot\text{ano}^{-1}$ (redução de $91{,}67\%$).
    \item $66{,}0$ para $11{,}0\text{ t}\cdot\text{ha}^{-1}\cdot\text{ano}^{-1}$ (redução de $83{,}33\%$).
    \item $110{,}0$ para $5{,}5\text{ t}\cdot\text{ha}^{-1}\cdot\text{ano}^{-1}$ (redução de $95{,}00\%$).
    \item $220{,}0$ para $11{,}0\text{ t}\cdot\text{ha}^{-1}\cdot\text{ano}^{-1}$ (redução de $95{,}00\%$).
    \item $33{,}0$ para $5{,}5\text{ t}\cdot\text{ha}^{-1}\cdot\text{ano}^{-1}$ (redução de $83{,}33\%$).
  \end{alternativas}

  \item \textbf{[Cap. 12, Protocolos de Campo]} Sobre os métodos quantitativos de monitoramento erosivo em campo, assinale a alternativa \textbf{INCORRETA}:
  \begin{alternativas}
    \item No método dos pinos de erosão, o aumento da altura exposta do pino ao longo do tempo indica rebaixamento erosivo do terreno.
    \item Perfis transversais repetidos em seções fixas de uma ravina permitem calcular o volume de solo mobilizado entre campanhas.
    \item Estacas externas posicionadas fora da borda de ruptura permitem monitorar a velocidade de recuo remontante da cabeceira de voçorocas.
    \item O fator $SDR$ (Razão de Aporte de Sedimentos) permite estimar que fração do solo erodido nas vertentes efetivamente chega ao reservatório ou exutório da bacia.
    \item A medição de apenas um pino isolado no topo da encosta dispensa réplicas espaciais e representa com exatidão o balanço de sedimentos de toda a bacia hidrográfica.
  \end{alternativas}
\end{enumerate}

\section*{Parte II: Questões Dissertativas Aplicadas}

\begin{enumerate}[leftmargin=1.8em, itemsep=8pt, start=9]
  \item \textbf{[Integração Caps. 10 e 11, PEIR e Qualidade do Solo no Semiárido]} Imagine que você foi contratado(a) para elaborar o diagnóstico ambiental dos acessos internos de um parque eólico recém-instalado sobre Caatinga na região de Canudos (BA).
  \begin{enumerate}[label=(\alph*), leftmargin=1.6em, itemsep=2pt]
    \item Estruture uma cadeia causal completa no modelo \textbf{PEIR} (Pressão, Estado, Impacto e Resposta) com pelo menos um indicador mensurável por componente.
    \item Selecione um \textbf{Conjunto Mínimo de Dados (MDS)} com $3$ indicadores de qualidade do solo (um físico, um químico e um biológico), explicando por que cada um deve ser comparado a uma área de referência na mesma classe de solo.
  \end{enumerate}
  \espacoResolucao[4.0cm]

  \item \textbf{[Integração Cap. 12, USLE, Rebaixamento e Erosão Linear]} Uma área de solo exposto ($C = 1{,}00$; $P = 1{,}00$) em uma vertente agrícola apresenta $R = 6.000\text{ MJ}\cdot\text{mm}\cdot\text{ha}^{-1}\cdot\text{h}^{-1}\cdot\text{ano}^{-1}$, $K = 0{,}03\text{ t}\cdot\text{ha}\cdot\text{h}\cdot\text{ha}^{-1}\cdot\text{MJ}^{-1}\cdot\text{mm}^{-1}$, $LS = 2{,}0$ e densidade do solo $D_s = 1{,}50\text{ g}\cdot\text{cm}^{-3}$. Além da erosão difusa, há uma ravina com $1{,}8\text{ m}$ de profundidade avançando sobre o terço médio da encosta.
  \begin{enumerate}[label=(\alph*), leftmargin=1.6em, itemsep=2pt]
    \item Calcule a perda anual de solo atual pela USLE ($A_1$) e a espessura correspondente de rebaixamento superficial ($h_1$, em $\text{mm}\cdot\text{ano}^{-1}$).
    \item Calcule a nova perda estimada ($A_2$) e a redução percentual caso se adote revegetação/manejo conservacionista com $C = 0{,}04$ e $P = 0{,}50$.
    \item Explique por que a USLE, isoladamente, não contabiliza o volume escavado na ravina de $1{,}8\text{ m}$ e indique qual protocolo de campo deve ser adotado para monitorá-la.
  \end{enumerate}
  \espacoResolucao[4.2cm]
\end{enumerate}

\clearpage
\section*{Gabarito Comentado e Espelho de Correção: Simulado do Módulo 4}

\begin{boxGabarito}[Gabarito Comentado: Parte I (Questões Objetivas 1 a 8)]
\begin{enumerate}[leftmargin=1.6em, itemsep=4pt]
  \item \textbf{(b)}, A troca de um protocolo financeiramente inviável por medições padronizadas e repetíveis atende diretamente ao critério de \textbf{mensurabilidade} (OECD, 2003; IBGE, 2018), assegurando continuidade da série histórica.
  \item \textbf{(b)}, Apenas (B) respeita o nexo causal correto: ação de obra (\textbf{P}) $\rightarrow$ alteração do solo (\textbf{E}) $\rightarrow$ danos erosivos e à fauna (\textbf{I}) $\rightarrow$ obras de controle, bioengenharia e ordenamento (\textbf{R}).
  \item \textbf{(c)}, O eixo Sobradinho--Paulo Afonso concentra múltiplas infraestruturas sobre o bioma Caatinga, exigindo indicadores integrados de água, solo, biota e conflitos territoriais.
  \item \textbf{(b)}, $D_s = 159 / 100 = 1{,}59\text{ g}\cdot\text{cm}^{-3}$; $P_t = (1 - 1{,}59 / 2{,}65)\times 100\% = (1 - 0{,}60)\times 100\% = \mathbf{40{,}0\%}$. Pelo fator de van Bemmelen: $CO = 17{,}24 / 1{,}724 = \mathbf{10{,}0\text{ g}\cdot\text{kg}^{-1}}$.
  \item \textbf{(c)}, Pontuações normalizadas: $S_{MO} = 24/30 = 0{,}80$; $S_{TIB} = 30/60 = 0{,}50$; $S_{D_s} = 1{,}20/1{,}50 = 0{,}80$ (menor melhor); $S_{DMP} = 1{,}50/3{,}00 = 0{,}50$. Média: $\text{IQS}_{\text{simples}} = (0{,}80 + 0{,}50 + 0{,}80 + 0{,}50)/4 = 2{,}60 / 4 = \mathbf{0{,}650}$.
  \item \textbf{(b)}, Crosta impermeável em solo nu = \textbf{selamento superficial}; canais rasos ($0{,}25\text{ m} < 0{,}5\text{ m}$) = \textbf{sulcos erosivos}; fenda profunda com surgência subterrânea e desabamento = \textbf{voçoroca ativa}.
  \item \textbf{(a)}, Potencial natural: $PNE = 5.500 \times 0{,}02 \times 2{,}0 = 220{,}0\text{ t}\cdot\text{ha}^{-1}\cdot\text{ano}^{-1}$. Cenário atual: $A_1 = 220{,}0 \times 0{,}30 \times 1{,}00 = \mathbf{66{,}0\text{ t}\cdot\text{ha}^{-1}\cdot\text{ano}^{-1}}$. Cenário conservacionista: $A_2 = 220{,}0 \times 0{,}05 \times 0{,}50 = \mathbf{5{,}5\text{ t}\cdot\text{ha}^{-1}\cdot\text{ano}^{-1}}$. Redução: $(66{,}0 - 5{,}5)/66{,}0 \times 100\% = 60{,}5/66{,}0 \times 100\% = \mathbf{91{,}67\%}$.
  \item \textbf{(E)}, Incorreta. A erosão apresenta elevada variabilidade espacial (zonas de perda e deposição); um único pino no topo não representa o balanço de sedimentos de toda a bacia, exigindo malha amostral distribuída.
\end{enumerate}
\end{boxGabarito}

\begin{boxGabarito}[Espelho de Correção: Parte II (Questões Dissertativas 9 e 10)]
\textbf{Questão 9 (2,0 pontos):}
\begin{itemize}[leftmargin=1.4em, itemsep=2.5pt]
  \item \textbf{Item (a), Cadeia PEIR (1,0 pt):} deve apresentar sequência lógica sem trocar papéis. Exemplo pleno: \textbf{P} = área de Caatinga suprimida ($\text{ha}$) e densidade viária ($D_v$, $\text{km}\cdot\text{km}^{-2}$); \textbf{E} = percentual de solo exposto ($\%$) e densidade do solo ($D_s$, $\text{g}\cdot\text{cm}^{-3}$); \textbf{I} = densidade de focos erosivos ativos ($\text{focos}\cdot\text{km}^{-1}$) e taxa corrigida de mortalidade de avifauna ($\text{aves}\cdot\text{torre}^{-1}\cdot\text{ano}^{-1}$); \textbf{R} = percentual de taludes revegetados em PRAD, obras de drenagem/bioengenharia executadas e monitoramento faunístico.
  \item \textbf{Item (b), MDS e área de referência (1,0 pt):} propor três indicadores complementares (por exemplo, físico como $D_s$, $P_{\text{macro}}$ ou $TIB$; químico como $CO$/$MO$, pH ou $V\%$; e biológico como $C\text{-BM}$, cobertura vegetal/serapilheira ou fauna edáfica) e justificar que a comparação com vegetação nativa sobre a \textbf{mesma classe pedológica e textura} isola a degradação antrópica das características inerentes do solo.
\end{itemize}

\textbf{Questão 10 (2,0 pontos):}
\begin{itemize}[leftmargin=1.4em, itemsep=2.5pt]
  \item \textbf{Item (a), Cálculo atual e rebaixamento (0,8 pt):}
  $A_1 = 6.000 \times 0{,}03 \times 2{,}0 \times 1{,}00 \times 1{,}00 = \mathbf{360{,}0\text{ t}\cdot\text{ha}^{-1}\cdot\text{ano}^{-1}}$.
  Rebaixamento: $h_1 = \frac{360{,}0}{10 \times 1{,}50} = \frac{360{,}0}{15{,}0} = \mathbf{24{,}0\text{ mm}\cdot\text{ano}^{-1}}$ ($2{,}4\text{ cm}\cdot\text{ano}^{-1}$).
  \item \textbf{Item (b), Novo cenário e redução percentual (0,6 pt):}
  $A_2 = 360{,}0 \times 0{,}04 \times 0{,}50 = \mathbf{7{,}2\text{ t}\cdot\text{ha}^{-1}\cdot\text{ano}^{-1}}$.
  Redução: $\frac{360{,}0 - 7{,}2}{360{,}0}\times 100\% = \frac{352{,}8}{360{,}0}\times 100\% = \mathbf{98{,}0\%}$.
  \item \textbf{Item (c), Limite da USLE e monitoramento da ravina (0,6 pt):} explicar que a USLE estima apenas erosão laminar e em sulcos rasos, não modelando erosão linear profunda em ravinas e voçorocas. Para monitorar a ravina de $1{,}8\text{ m}$, devem-se empregar \textbf{perfis transversais periódicos} em seções fixas (cubagem de volume $V$ e massa $M$) e \textbf{estacas de referência externas} para medir o recuo remontante da cabeceira, além de ficha visual para verificar eventual surgimento de \textit{piping}.
\end{itemize}
\end{boxGabarito}

\section*{Tabela de Diagnóstico de Revisão do Módulo 4}

Registre seu desempenho no simulado para direcionar a revisão antes de avançar para o Módulo 5 (Bioengenharia de Água e Solo, Legislação Ambiental e PRAD):

\begin{center}
\renewcommand{\arraystretch}{1.35}
\small
\begin{tabularx}{\textwidth}{@{}p{2.2cm} p{4.2cm} c c >{\raggedright\arraybackslash}X@{}}
  \toprule
  \rowcolor{uefsNavy}
  \textcolor{white}{\bfseries Capítulo} & \textcolor{white}{\bfseries Eixo Avaliado} & \textcolor{white}{\bfseries Questões} & \textcolor{white}{\bfseries Acertos} & \textcolor{white}{\bfseries Ação Recomendada se Abaixo de 70\%} \\
  \midrule
  \textbf{Cap. 10} & Critérios OCDE, PEIR e Campo no Semiárido & Q1, Q2, Q3, Q9a & \quad / 4 & Refazer Exemplos 10.1 a 10.4 e revisar a Figura~\ref{fig:modelo_peir_semiarido}. \\
  \rowcolor{bgGray}
  \textbf{Cap. 11} & Física/Química do Solo, van Bemmelen e IQS & Q4, Q5, Q9b & \quad / 3 & Refazer Exemplos 11.1, 11.2 e 11.5 e revisar normalização (``menor melhor''). \\
  \textbf{Cap. 12} & Erosão Linear, USLE, Pinos e Perfis & Q6, Q7, Q8, Q10 & \quad / 4 & Refazer Exemplos 12.1 a 12.4 e revisar a Figura~\ref{fig:evolucao_erosao_linear}. \\
  \bottomrule
\end{tabularx}
\end{center}
